\documentclass[11pt]{article}
\usepackage[T1]{fontenc}
\usepackage{lmodern}
\usepackage[margin=1in]{geometry}
\usepackage{amsmath,amssymb,amsthm,mathtools}
\usepackage{microtype}
\usepackage{booktabs,longtable,array}
\usepackage{enumitem}
\usepackage{xcolor}
\usepackage[hidelinks]{hyperref}
\usepackage{fancyhdr}
\setlength{\headheight}{14pt}
\pagestyle{fancy}
\fancyhf{}\fancyhead[L]{\small Cubic invariant tail descent}\fancyhead[R]{\small ProveIt research continuation}\fancyfoot[C]{\thepage}
\setlength{\parskip}{4pt}
\setlength{\parindent}{0pt}
\allowdisplaybreaks[2]
\numberwithin{equation}{section}
\newtheorem{theorem}{Theorem}[section]
\newtheorem{proposition}[theorem]{Proposition}
\newtheorem{lemma}[theorem]{Lemma}
\newtheorem{corollary}[theorem]{Corollary}
\theoremstyle{definition}\newtheorem{definition}[theorem]{Definition}
\theoremstyle{remark}\newtheorem{remark}[theorem]{Remark}
\DeclareMathOperator{\Li}{Li}
\DeclareMathOperator{\Rea}{Re}
\DeclareMathOperator{\depth}{depth}
\newcommand{\dd}{\,\mathrm d}
\newcommand{\QQ}{\mathbb Q}
\newcommand{\CC}{\mathbb C}
\newcommand{\NN}{\mathbb N}
\newcommand{\Om}{\Omega}
\newcommand{\Ztr}{\mathcal Z}
\newcommand{\Ltr}{\mathcal L}
\newcommand{\Dtr}{\mathcal D}
\newcommand{\coef}[1]{\left[#1\right]}
\newcommand{\src}[1]{\nolinkurl{#1}}
\hypersetup{pdftitle={Cubic Invariant Tail Descent},pdfauthor={Research continuation prepared for Vladimir Reshetnikov},pdfsubject={Exact harmonic-tail identities, Gamma quotients, polylogarithms and Stieltjes traces}}
\title{\textbf{Cubic Invariant Tail Descent}\\[5pt]
\large All-degree harmonic identities, exact digamma cubes,\\
polylogarithmic representations, and algebraic Stieltjes traces}
\author{Research continuation prepared for Vladimir Reshetnikov\\
\small ProveIt / Analysis / Polylogarithms\\
\small Mathematical derivations and computational checks with OpenAI assistance}
\date{10 October 2026}
\begin{document}
\maketitle
\begin{abstract}
The first mixed Rogers--Dougall resonance already present in the ProveIt
manuscript has a symmetric cubic structure that is not visible in its original
two-parameter coefficient tables. We expose that structure and prove an
all-degree descent theorem: the weighted, symmetrized product of $d\ge2$
odd zeta tails has an explicit expression in multiple zeta values of depth
at most $\max(1,d-2)$. The weight is preserved after the natural shift by
two. A finite recurrence supplies every coefficient.

The first consequences include
\[
 \sum_{n\ge0}(2n+1)\bigl(\zeta(3)-H_n^{(3)}\bigr)^3
       =\frac74\zeta(7),
\]
and
\[
 \sum_{n\ge0}(2n+1)\bigl(\zeta(3)-H_n^{(3)}\bigr)^4
       =\zeta(10)+\zeta(5)^2+2\zeta(7,3).
\]
We also obtain meromorphic closed formulas for the square and cube of a
symmetric digamma tail, an exact cubic identity involving $H_n,H_{2n}$ and
$\log2$, and convergent polylogarithmic integral and primitive formulas.
The same cubic coordinates give a branch-independent expansion of Hurwitz
and Lerch traces, all-index generalized Stieltjes identities, and anchored
antiderivatives of every positive integer order. Root collisions cause no
singularity in these trace formulas.

The proofs are analytic and algebraic, not integer-relation inferences.
The accompanying finite exact checks and numerical diagnostics are separate
from the proofs. The classical Gamma and tail-product inputs are attributed;
no comprehensive historical-priority claim is made. The Gaussian $S_6$ and
revised $S_8$ conjectures remain unresolved by this article.
\end{abstract}
\vfill
\noindent\small\textbf{Source scope.} Targeted canonical reads and repository
searches at the source anchor\\
\src{87c7db629552afbc2e803001526debe3430a31b0}.
The incoming directory inventory and intake instructions were inspected;
the binary contents of its ZIPs were not audited. No live repository files
were modified.\normalsize
\clearpage
\begingroup\small\setlength{\parskip}{0pt}
\tableofcontents
\endgroup
\clearpage

\section{Research position and principal results}
\label{ct:sec:position}
The relevant canonical starting point is
\src{chapters/07-mixed-dougall-identities.tex} \cite{ct:repo-mixed}.
It proves a normally convergent first-resonance Gamma identity, mixed
parameter derivatives, a weighted trigamma-difference square, and an
all-order quadratic zeta-tail reduction. Those results are inputs here,
not new claims. The harmonic coefficient chapter separately proves the
Gauss identities that resolve and correct the indicated Choi formulas
\cite{ct:repo-harmonic}. We do not reopen those resolved conjectures.

The contribution of this continuation is a precise structural consequence
of the mixed resonance. Introducing the elementary invariants of the
three roots $p,q,-p-q$ turns a large family of mixed derivatives into a
triangular system ordered by the number of tail factors. Its leading
term is an arbitrary symmetric convolution of odd tails; every other term
has at least two fewer factors. This gives an explicit all-degree depth
reduction, rather than another unorganized list of derivatives.

Let
\begin{equation}\label{ct:eq:tails}
 T_s(n)=\zeta(s,n+1)=\zeta(s)-H_n^{(s)},\qquad
 H_n^{(s)}=\sum_{j=1}^n j^{-s},\qquad n\ge0,
\end{equation}
initially with integer $s\ge3$. In particular $H_0^{(s)}=0$.
For $a\ge0$ and $d\ge2$, set
\begin{equation}\label{ct:eq:M}
 M_{a,d}=\sum_{n\ge0}(2n+1)
       \sum_{j_1+\cdots+j_d=a\atop j_i\ge0}
             \prod_{i=1}^d T_{2j_i+3}(n).
\end{equation}
The index tuples in the inner sum are ordered. This multiplicity is
important: $M_{1,3}=3\sum_n(2n+1)T_3(n)^2T_5(n)$.

\begin{theorem}[Summary of the tail descent]\label{ct:thm:summary}
For every $a\ge0,d\ge2$, $M_{a,d}$ has an explicit rational linear
expression in convergent multiple zeta values of weight
$2a+3d-2$ and depth at most $\max(1,d-2)$. The construction is the finite
recurrence in Theorem~\ref{ct:thm:coefficient} followed by the ordered-partition
map in Lemma~\ref{ct:lem:partitions}. In particular,
\begin{align}
 M_{a,2}&=\frac{2a+3}{2}\zeta(2a+4),\label{ct:eq:quadratic}\\
 M_{a,3}&=\frac{(a+1)(2a+7)}4\zeta(2a+7).\label{ct:eq:cubic}
\end{align}
\end{theorem}
Here and below ``depth at most'' is a constructive upper bound on a
representation. It is not a minimal-depth assertion, a statement of period
independence, or a proof that $\zeta(7,3)$ cannot be reduced further.

For a second accessible family, define
\begin{equation}\label{ct:eq:R}
 R_n=2\bigl(\log2+H_n-H_{2n}\bigr)-\frac1{2n+1}.
\end{equation}
We prove
\begin{equation}\label{ct:eq:R-identities}
 \sum_{n\ge0}(2n+1)R_n^2=\frac{\zeta(2)-1}{4},\qquad
 \sum_{n\ge0}(2n+1)R_n^3=\frac58\zeta(3)-\log2.
\end{equation}
These are ordinary absolutely convergent series. Their proof gives a full
complex-parameter identity, not just these two evaluations.

\subsection{Classical antecedents and novelty boundaries}
The Rogers--Dougall very well-poised summation is classical
\cite{ct:dlmf16}. The passage from products of zeta tails to multiple zeta
values is also established; Hoffman's treatment is a direct antecedent
\cite{ct:hoffman}. Our minimum-square version of that passage is proved
below because the weight $2n+1$ and the inclusion of $n=0$ change the
normalization. It is not presented as the discovery of the general
sums-of-tails method.

The symmetric Gamma product, the Hurwitz shift expansion, and the
Gamma--zeta derivative at zero are standard analytic ingredients
\cite{ct:dlmf5,ct:dlmf25}. The additional results are their explicit cubic
invariant organization, the associated all-degree descent, the closed
symmetric-digamma cube, and the compatible trace/primitive formulas.
These statements are proved here. A targeted search did not establish
comprehensive historical priority for their individual special cases.

\section{The cubic Gamma quotient}
\label{ct:sec:gamma}
Let
\begin{equation}\label{ct:eq:domain}
 \Om=\{(\sigma,\tau)\in\CC^2:|\sigma|+|\tau|<1\}.
\end{equation}
Write $r_1,r_2,r_3$ for the roots, with multiplicity, of
\begin{equation}\label{ct:eq:cubic-poly}
 r^3-\sigma r+\tau=0.
\end{equation}
Rouch\'e's theorem on $|r|=1$ gives $|r_j|<1$ on $\Om$.
The symmetric expressions below do not require a continuous labeling of
these roots. Put
\begin{align}
 Q_n(\sigma,\tau)
   &=\prod_{j=1}^3\frac{\Gamma(n+1+r_j)}{\Gamma(n+1-r_j)},
       \qquad x_n=n+\tfrac12,\label{ct:eq:Q}\\
 B(\sigma,\tau)
   &=\psi(\tfrac12)+\frac12\left\{\psi(2)-\sum_{j=1}^3\psi(1+r_j)\right\}.
       \label{ct:eq:B}
\end{align}

\begin{lemma}[Product and logarithm]\label{ct:lem:product}
On $\Om$,
\begin{equation}\label{ct:eq:product}
 Q_n(\sigma,\tau)
   =\prod_{j>n}\frac{j^3-\sigma j+\tau}{j^3-\sigma j-\tau},
\end{equation}
and the logarithm that is zero at $(0,0)$ has the locally normally
convergent expansion
\begin{equation}\label{ct:eq:logQ}
 \log Q_n
 =2\sum_{a,b\ge0}\frac{\binom{a+2b}{a}}{2b+1}
       T_{2a+6b+3}(n)\,\sigma^a\tau^{2b+1}.
\end{equation}
Moreover,
\begin{equation}\label{ct:eq:Q-bound}
 x_n(Q_n-1)-\frac{\tau}{x_n}=O(n^{-3})
\end{equation}
uniformly on compact subsets of $\Om$.
\end{lemma}
\begin{proof}
The relations among the roots are
$\sum r_j=0$, $\sum_{i<j}r_ir_j=-\sigma$ and $r_1r_2r_3=-\tau$.
The linear exponential factors in the Weierstrass product for Gamma cancel.
The remaining product gives \eqref{ct:eq:product}. Equivalently, use the
Gamma recurrence to remove its first $n$ factors. Both numerator and
denominator in \eqref{ct:eq:product} are nonzero on $\Om$.

For $j\ge1$, expand
\[
 \log\frac{1+\tau/(j^3-\sigma j)}{1-\tau/(j^3-\sigma j)}
   =2\sum_{b\ge0}\frac{\tau^{2b+1}}{(2b+1)(j^3-\sigma j)^{2b+1}},
\]
then expand $(1-\sigma/j^2)^{-2b-1}$. On every closed smaller domain
$|\sigma|+|\tau|\le\rho<1$, these geometric expansions admit an absolute
majorant. Summation over $j>n$ proves \eqref{ct:eq:logQ} and its normal
convergence.

The first term is $2\tau T_3(n)$. All other terms are $O(n^{-4})$ uniformly
on compact parameter sets. The elementary centered tail expansion is
\[
 T_3(n)=\frac1{2x_n^2}+O(x_n^{-4}).
\]
For example, it follows from the first two terms of Euler--Maclaurin in
$n+1$ followed by the substitution $n+1=x_n+1/2$; the cubic inverse power
cancels. Thus $\log Q_n=\tau/x_n^2+O(x_n^{-4})$ and exponentiation proves
\eqref{ct:eq:Q-bound}.
\end{proof}

\begin{theorem}[Cubic invariant master identity]\label{ct:thm:master}
For $(\sigma,\tau)\in\Om$,
\begin{equation}\label{ct:eq:master}
 \sum_{n\ge0}\left\{x_n(Q_n(\sigma,\tau)-1)-\frac{\tau}{x_n}\right\}
       =\tau B(\sigma,\tau).
\end{equation}
The bracketed sum is locally normally convergent. Every fixed derivative
in $\sigma,\tau$ may be taken termwise, provided the displayed subtraction
is retained before differentiation.
\end{theorem}

This theorem is the $a=1$ specialization of the canonical mixed resonance,
rewritten symmetrically. To make the present continuation independently
readable, the next section supplies a direct derivation from the classical
summation, including the finite constant. It is not a second novelty claim
for that resonance.

\section{A proof of the master identity with its subtraction visible}
\label{ct:sec:master-proof}
Initially choose small negative real $p,q$ and then a sufficiently small
complex neighborhood of a generic such pair. Put $d_0=-p-q$, with
$\Rea d_0>0$, and set
\[
 \sigma=p^2+pq+q^2,\qquad \tau=pq(p+q).
\]
Then the three roots are $p,q,-p-q$. Introduce a separate spectral parameter
$u$ and define
\begin{equation}\label{ct:eq:ru}
 r_n(u)=x_n\frac{\Gamma(n+1+p)\Gamma(n+1+q)\Gamma(n-p-q-u)}
 {\Gamma(n+1-p)\Gamma(n+1-q)\Gamma(n+2+p+q+u)}.
\end{equation}
For $0<\Rea u<\Rea d_0$, the Rogers--Dougall summation, in the normalization of
\eqref{ct:eq:ru}, is
\begin{equation}\label{ct:eq:G}
 \sum_{n\ge0}r_n(u)=G(u):=
 \frac{\Gamma(1+p)\Gamma(1+q)\Gamma(-p-q-u)\Gamma(u)}
 {2\Gamma(-p-q)\Gamma(1+p+u)\Gamma(1+q+u)}.
\end{equation}
Indeed substitute the upper parameters
$1,3/2,1+p,1+q,-p-q-u$ in the classical ${}_5F_4(1)$ formula
\cite[Eq.~16.4.9]{ct:dlmf16}; its excess is $2u$.
The factor $1/2$ comes from $x_n=n+1/2$.

The Gamma ratio in \eqref{ct:eq:ru} is centered at $x_n$. The numerator
shifts, including the canceled pair $\Gamma(n+1)/\Gamma(n+1)$, are
\[
 \alpha=(\tfrac12,\tfrac12+p,\tfrac12+q,-\tfrac12-p-q-u),
\]
and the denominator shifts are $1-\alpha_j$. The shifted logarithmic
Gamma expansion and $B_k(1-z)=(-1)^kB_k(z)$ give
\begin{equation}\label{ct:eq:asymptotic-r}
 r_n(u)=x_n^{-1-2u}\{1+C_1(u)x_n^{-2}+O(x_n^{-4})\},
 \qquad C_1(u)=-\frac13\sum_{j=1}^4B_3(\alpha_j).
\end{equation}
The odd inverse powers cancel before exponentiation. These estimates are
uniform on compact parameter sets avoiding the displayed Gamma poles
\cite{ct:dlmf5}. Consequently
\begin{align}
 \mathcal R(u)&=\sum_{n\ge0}
 \{r_n(u)-x_n^{-1-2u}-C_1(u)x_n^{-3-2u}\}\\
 &=G(u)-\zeta(1+2u,\tfrac12)-C_1(u)\zeta(3+2u,\tfrac12)
 \label{ct:eq:R-cont}
\end{align}
continues holomorphically across $u=-1$. Locally the summands are
$O(n^{-5-2\Rea u})$, so the continuation is also the normally convergent
subtracted series there.

Write $v=u+1$. Direct Bernoulli algebra gives
\begin{equation}\label{ct:eq:C1}
 C_1(-1)=\tau,\qquad
 \frac12C_1'(-1)+\zeta(-1,\tfrac12)=\frac{(p+q)^2}{2}.
\end{equation}
For clarity, $B_3(1/2+t)=t^3-t/4$, and
$C_1'(-1)=B_2(1/2-p-q)=(p+q)^2-1/12$;
$\zeta(-1,1/2)=1/24$.

The Laurent constant of $G(-1+v)$ is
\begin{equation}\label{ct:eq:G-constant}
 G(-1+v)=\frac{\tau}{2v}
 +\frac\tau2\{\psi(2)-\psi(1-p-q)-\psi(p)-\psi(q)\}+O(v).
\end{equation}
The resonant Hurwitz term is
\[
 C_1(-1+v)\zeta(1+2v,\tfrac12)
 =\frac\tau{2v}+\frac12C_1'(-1)-\tau\psi(\tfrac12)+O(v).
\]
Subtracting these expansions in \eqref{ct:eq:R-cont}, and using
$\psi(p)=\psi(1+p)-1/p$ and its $q$ counterpart, cancels precisely
\[
 \frac\tau2\left(\frac1p+\frac1q\right)
       =\frac{(p+q)^2}{2}.
\]
Since $r_n(-1)=x_nQ_n$, the result is \eqref{ct:eq:master}.

The initial restrictions on $p,q$ can now be removed by holomorphy. Both
sides are symmetric in the three roots. Alternatively, regard them as
holomorphic functions of $\sigma,\tau$ using the product and power-series
representations. The map $(p,q)\mapsto(\sigma,\tau)$ has nonzero Jacobian
at generic small pairs, so the identity holds on a nonempty open subset of
$\Om$. The identity theorem and connectedness of $\Om$ extend it throughout.
Lemma~\ref{ct:lem:product} gives normal convergence directly on this entire
domain. Cauchy's integral formula on a slightly larger compact polydisc
then justifies all fixed coefficient extractions and derivatives.
\hfill$\square$

\section{Newton coefficients and the all-degree tail recurrence}
\label{ct:sec:coefficients}
Put $S_k=r_1^k+r_2^k+r_3^k$. The relation
\begin{equation}\label{ct:eq:newton-log}
 -\log(1-\sigma t^2+\tau t^3)=\sum_{k\ge1}\frac{S_k}{k}t^k
\end{equation}
gives, for $k\ge1$,
\begin{equation}\label{ct:eq:newton}
 S_k=\sum_{2a+3b=k\atop a,b\ge0}
 \frac{k}{a+b}\binom{a+b}{b}(-1)^b\sigma^a\tau^b.
\end{equation}
There is no $(a,b)=(0,0)$ term. Equivalently,
$S_0=3,S_1=0,S_2=2\sigma$ and $S_k=\sigma S_{k-2}-\tau S_{k-3}$.

Using $\psi(1+z)=-\gamma+\sum_{k\ge1}(-1)^{k+1}\zeta(k+1)z^k$,
the two signs $(-1)^k$ and $(-1)^b$ cancel when $k=2a+3b$. Hence
\begin{equation}\label{ct:eq:B-series}
 B(\sigma,\tau)=\frac12-2\log2+
 \frac12\sum_{a,b\ge0\atop a+b>0}
 \frac{2a+3b}{a+b}\binom{a+b}{b}
       \zeta(2a+3b+1)\sigma^a\tau^b.
\end{equation}
The expansion is normally convergent on smaller polydiscs inside $\Om$.
The power-sum formula, or the root trace itself, continues it wherever
needed in $\Om$.

\begin{definition}\label{ct:def:P}
Treat $T_3,T_5,\ldots$ as independent formal symbols. Define
\begin{equation}\label{ct:eq:P-def}
 P_{a,d}=[\sigma^a\tau^d]\exp\left(
 2\sum_{i,b\ge0}\frac{\binom{i+2b}{i}}{2b+1}
       T_{2i+6b+3}\sigma^i\tau^{2b+1}\right).
\end{equation}
This is a finite polynomial for each fixed $(a,d)$.
\end{definition}

\begin{theorem}[Finite recurrence and exact summation]\label{ct:thm:coefficient}
The initial values are $P_{0,0}=1$ and $P_{a,0}=0$ for $a>0$. For $d\ge1$,
\begin{equation}\label{ct:eq:recurrence}
 P_{a,d}=\frac2d\sum_{b\le d\atop b\text{ odd}}
            \sum_{i=0}^a\binom{i+b-1}{i}
                T_{2i+3b}\,P_{a-i,d-b}.
\end{equation}
For every $a\ge0,d\ge2$,
\begin{equation}\label{ct:eq:P-sum}
 \sum_{n\ge0}(2n+1)P_{a,d}\bigl(T_3(n),T_5(n),\ldots\bigr)
       =C_{a,d}\zeta(2a+3d-2),
\end{equation}
where
\begin{equation}\label{ct:eq:C-ad}
 C_{a,d}=\frac{2a+3d-3}{a+d-1}\binom{a+d-1}{d-1}.
\end{equation}
All the sums in \eqref{ct:eq:P-sum} converge absolutely.
\end{theorem}
\begin{proof}
Differentiating the exponential in \eqref{ct:eq:P-def} with respect to
$\tau$ gives \eqref{ct:eq:recurrence}. On substituting the genuine tails,
$P_{a,d}$ is exactly $[\sigma^a\tau^d]Q_n$. Take this coefficient in
Theorem~\ref{ct:thm:master} and multiply by two. When $d\ge2$, the
counterterm and the constant term of $B$ make no contribution.
The coefficient in \eqref{ct:eq:B-series} with $b=d-1$ gives
\eqref{ct:eq:C-ad}.

Each monomial is a product of $\ell\ge1$ tails, all of orders at least
three, with total order $2a+3d$. Because $T_s(n)=O(n^{1-s})$, its weighted
summand is $O(n^{1+\ell-2a-3d})$. Here $\ell\le d$, so the exponent is
at most $1-2a-2d\le-3$. This also directly proves absolute convergence
without relying only on coefficient extraction.
\end{proof}

\subsection{The exceptional linear coefficient}
The sole coefficient requiring the subtraction is $(a,d)=(0,1)$:
\begin{equation}\label{ct:eq:linear-counterterm}
 \sum_{n\ge0}\left\{(2n+1)T_3(n)-\frac1{n+1/2}\right\}
       =\frac12-2\log2.
\end{equation}
For $a\ge1,d=1$, the identity is simply
$\sum_n(2n+1)T_{2a+3}(n)=\zeta(2a+1)$.
Equation~\eqref{ct:eq:linear-counterterm} must not be split into the
difference of two ordinarily convergent sums: both separate series diverge.
No such finite-part issue is present in the nonlinear identities.

\section{Why two levels of depth disappear}
\label{ct:sec:descent}
Use the convention
\begin{equation}\label{ct:eq:mzv}
 \zeta(s_1,\ldots,s_\ell)
  =\sum_{m_1>\cdots>m_\ell\ge1}
          m_1^{-s_1}\cdots m_\ell^{-s_\ell}.
\end{equation}
An ordered set partition of $\{1,\ldots,d\}$ is a list of nonempty disjoint
blocks whose union is that set; the order of the blocks matters.

\begin{lemma}[The minimum-square map]\label{ct:lem:partitions}
For integers $s_i\ge3$,
\begin{equation}\label{ct:eq:partition-map}
 \sum_{n\ge0}(2n+1)\prod_{i=1}^dT_{s_i}(n)
  =\sum_{(B_1,\ldots,B_\ell)}
     \zeta\left(\sum_{i\in B_1}s_i,\ldots,
                  \sum_{i\in B_\ell}s_i-2\right),
\end{equation}
where the last exponent alone is reduced by two. A one-block partition
contributes $\zeta(s_1+\cdots+s_d-2)$.
\end{lemma}
\begin{proof}
Expand the product as a sum over $m_i>n$. Positivity permits Tonelli, and
\[
 \sum_{0\le n<\min(m_1,\ldots,m_d)}(2n+1)
       =\min(m_1,\ldots,m_d)^2.
\]
Group the indices $m_i$ according to their equal values, listing the blocks
from largest value to smallest. The minimum-square factor reduces the
last exponent by two and produces exactly the summand in
\eqref{ct:eq:partition-map}. No factorial is inserted: every ordered
partition already records one distinct equality-and-order pattern.
All resulting multiple zeta series converge. Their first exponent is at
least three unless there is just one block, and their final exponent is
at least one. Alternatively, the preceding tail estimate proves the
finiteness of the original positive sum.
\end{proof}
This is the weighted counterpart of the classical tail-product
rearrangement discussed by Hoffman \cite{ct:hoffman}.

Define
\[
 A_n(\sigma)=\sum_{j\ge0}T_{2j+3}(n)\sigma^j.
\]
The part of $\log Q_n$ linear in $\tau$ is $2\tau A_n(\sigma)$.
Every remaining term has odd $\tau$-degree at least three. Consequently
\begin{equation}\label{ct:eq:leading}
 P_{a,d}=\frac{2^d}{d!}[\sigma^a]
          \left(\sum_{j\ge0}T_{2j+3}\sigma^j\right)^d+E_{a,d},
\end{equation}
where every monomial of $E_{a,d}$ has at most $d-2$ tail factors and the
same parity of factor count as $d$. This is the entire reason for the
depth drop; it is not a numerical rank observation.

\begin{theorem}[Constructive depth descent]\label{ct:thm:descent}
For $a\ge0,d\ge2$,
\begin{equation}\label{ct:eq:descent}
 M_{a,d}=\frac{d!}{2^d}\left\{
 C_{a,d}\zeta(2a+3d-2)
 -\sum_{n\ge0}(2n+1)E_{a,d}\bigl(T_3(n),T_5(n),\ldots\bigr)
 \right\}.
\end{equation}
Applying \eqref{ct:eq:partition-map} to the finitely many monomials in
$E_{a,d}$ produces only multiple zeta values of weight $2a+3d-2$ and depth
at most $d-2$. For $d=2$, $E_{a,2}=0$, so the result has depth one.
\end{theorem}
\begin{proof}
The unique way to obtain $d$ tail factors at $\tau$-degree $d$ in the
exponential is to use $d$ linear terms. If any factor has $\tau$-degree
at least three, the number of factors decreases by at least two. This
proves \eqref{ct:eq:leading}. Substitution in \eqref{ct:eq:P-sum} proves
\eqref{ct:eq:descent}. Each monomial has total tail order $2a+3d$;
Lemma~\ref{ct:lem:partitions} reduces that weight by exactly two and never
increases its number of factors. All operations involve finite linear
combinations of absolutely convergent series.
\end{proof}
The theorem concerns the symmetrized convolution $M_{a,d}$, not every
individual product with the same weight. For instance it directly reduces
the sum of $T_3^2T_7$ and $T_3T_5^2$, but does not separately evaluate both.
That distinction prevents an unjustified polarization claim.

\section{Cubes, fourth powers, and a fifth-power relation}
\label{ct:sec:examples}
\subsection{The complete quadratic and cubic families}
The $\tau^2$ coefficient is $2A_n(\sigma)^2$, giving
\eqref{ct:eq:quadratic}. At degree three, put
\[
 C_n(\sigma)=\sum_{a\ge0}\binom{a+2}{2}T_{2a+9}(n)\sigma^a.
\]
Then
\begin{equation}\label{ct:eq:tau3}
 [\tau^3]Q_n=\frac43A_n(\sigma)^3+\frac23C_n(\sigma).
\end{equation}
The elementary identity $\sum_n(2n+1)T_s(n)=\zeta(s-2)$ for $s>3$ implies
\[
 \frac43M_{a,3}+\frac23\binom{a+2}{2}\zeta(2a+7)
       =(a+1)(a+3)\zeta(2a+7).
\]
Solving proves \eqref{ct:eq:cubic}. Explicit consequences are
\begin{align}
 \sum_n(2n+1)T_3(n)^3&=\frac74\zeta(7),\label{ct:eq:cube0}\\
 \sum_n(2n+1)T_3(n)^2T_5(n)&=\frac32\zeta(9),\label{ct:eq:cube1}\\
 \sum_n(2n+1)\{T_3(n)^2T_7(n)+T_3(n)T_5(n)^2\}
       &=\frac{11}{4}\zeta(11).\label{ct:eq:cube2}
\end{align}
For example, the first identity also states
\begin{equation}\label{ct:eq:cubic-mzv}
 6\zeta(3,3,1)+3\zeta(3,4)+3\zeta(6,1)+\zeta(7)
      =\frac74\zeta(7).
\end{equation}
This is an equality of periods with a proof from the master identity;
it is not an assertion that the four displayed coordinates form a basis.

\subsection{A fourth power reduces to one depth-two coordinate}
At $a=0$,
\begin{equation}\label{ct:eq:tau4}
 P_{0,4}=\frac23T_3^4+\frac43T_3T_9.
\end{equation}
Since $C_{0,4}=3$, Theorem~\ref{ct:thm:coefficient} gives
\begin{equation}\label{ct:eq:quartic-start}
 \sum_n(2n+1)T_3(n)^4
       =\frac92\zeta(10)-2\sum_n(2n+1)T_3(n)T_9(n).
\end{equation}
The minimum-square map yields
\[
 \sum_n(2n+1)T_3(n)T_9(n)
       =\zeta(10)+\zeta(3,7)+\zeta(9,1).
\]
Euler's depth-two reduction and the ordinary stuffle product give
\begin{align*}
 \zeta(9,1)&=\frac92\zeta(10)-\zeta(2)\zeta(8)-\zeta(3)\zeta(7)
                         -\zeta(4)\zeta(6)-\tfrac12\zeta(5)^2,\\
 \zeta(3,7)+\zeta(7,3)&=\zeta(3)\zeta(7)-\zeta(10).
\end{align*}
These classical relations are recorded, with the same series convention,
in \cite[\S2]{ct:hoffman}. Substituting them and using
$\zeta(2)\zeta(8)=33\zeta(10)/20$ and
$\zeta(4)\zeta(6)=11\zeta(10)/10$ proves the following.

\begin{corollary}\label{ct:cor:quartic}
\begin{equation}\label{ct:eq:quartic}
 \sum_{n\ge0}(2n+1)\bigl(\zeta(3)-H_n^{(3)}\bigr)^4
       =\zeta(10)+\zeta(5)^2+2\zeta(7,3).
\end{equation}
\end{corollary}

\subsection{The next descent, without claiming a complete basis reduction}
The fifth coefficient is
\[
 P_{0,5}=\frac4{15}T_3^5+\frac43T_3^2T_9+\frac25T_{15}.
\]
It follows that
\begin{equation}\label{ct:eq:quintic}
 \sum_n(2n+1)\{T_3(n)^5+5T_3(n)^2T_9(n)\}
       =\frac{39}{4}\zeta(13).
\end{equation}
For a fully explicit depth-three form, the remaining mixed moment is
\begin{align}
 \sum_n(2n+1)T_3(n)^2T_9(n)
   ={}&2\zeta(3,3,7)+2\zeta(3,9,1)+2\zeta(9,3,1)\notag\\
     &+\zeta(6,7)+\zeta(9,4)
           +2\zeta(12,1)+2\zeta(3,10)+\zeta(13).
 \label{ct:eq:339}
\end{align}
The coefficients count the ordered partitions of the two labeled $3$-tails
and one labeled $9$-tail. Thus \eqref{ct:eq:quintic} really is a finite
reduction of the fifth power to depth at most three. No numerical basis
selection is needed to make it a proved exact identity.

\section{Closed meromorphic formulas for symmetric digamma powers}
\label{ct:sec:digamma}
For $x\notin\mathbb Z\setminus\{0\}$, define
\begin{equation}\label{ct:eq:A-x}
 \mathcal A_n(x)=\sum_{j>n}\frac1{j(j^2-x^2)}
 =\frac{2\psi(n+1)-\psi(n+1-x)-\psi(n+1+x)}{2x^2},
\end{equation}
with $\mathcal A_n(0)=T_3(n)$. The equality follows by partial fractions
and the digamma difference series. For $|x|<1$,
$\mathcal A_n(x)=A_n(x^2)$.

Set
\begin{equation}\label{ct:eq:J}
 J(x)=\psi(1+x)+\psi(1-x)-2\psi(1).
\end{equation}
Then $J'(x)=\psi'(1+x)-\psi'(1-x)$ and
$J''(x)=\psi''(1+x)+\psi''(1-x)$.

\begin{theorem}[Exact square and cube]\label{ct:thm:digamma}
On $\CC\setminus(\mathbb Z\setminus\{0\})$,
\begin{align}
 \sum_{n\ge0}(2n+1)\mathcal A_n(x)^2
    &=\frac{\pi^2}{4x^2\sin^2(\pi x)}
        -\frac{\zeta(2)}{2x^2}-\frac1{4x^4},
       \label{ct:eq:A-square}\\
 \sum_{n\ge0}(2n+1)\mathcal A_n(x)^3
    &=\frac{J(x)}{4x^6}+\frac{J'(x)}{8x^5}
        -\frac{J''(x)}{16x^4}+\frac{3\zeta(3)}{4x^4}.
       \label{ct:eq:A-cube}
\end{align}
All singularities at $x=0$ on the right are removable. The powers are
algebraic powers, not absolute-value powers. The series converge locally
normally away from the nonzero integers.
\end{theorem}
\begin{proof}
First let $\sigma=x^2$, $|x|<1$, and $D=\sigma\partial_\sigma$.
The quadratic family gives
\[
 \sum_n(2n+1)\mathcal A_n(x)^2=(D+\tfrac32)F(\sigma),\qquad
 F(\sigma)=\sum_{a\ge0}\zeta(2a+4)\sigma^a.
\]
Using the partial-fraction expansion of cotangent,
\[
 E(\sigma)=\sum_{j\ge1}\frac1{j^2-\sigma}
       =\frac{1-\pi x\cot(\pi x)}{2x^2},\qquad
 F(\sigma)=\frac{E(\sigma)-\zeta(2)}{\sigma}.
\]
Differentiation cancels the cotangent terms and gives \eqref{ct:eq:A-square}.

Similarly, the cubic family yields
\begin{align*}
 \sum_n(2n+1)\mathcal A_n(x)^3
       &=\tfrac14(D+1)(2D+7)G(\sigma),\\
 G(\sigma)&=\sum_{a\ge0}\zeta(2a+7)\sigma^a
       =\frac{\mathcal A_0(x)-\zeta(3)-x^2\zeta(5)}{x^4}.
\end{align*}
Since $\mathcal A_0(x)=-J(x)/(2x^2)$ and $D=(x/2)\partial_x$, this operator
is $x^2\partial_x^2/8+5x\partial_x/4+7/4$.
Its action on $-J(x)/(2x^6)-\zeta(3)/x^4-\zeta(5)/x^2$ gives
\eqref{ct:eq:A-cube}; the $\zeta(5)$ term cancels identically.

On a compact set avoiding the nonzero integers, the finitely many small
$j$ denominators are harmless and the large-$n$ tail satisfies
$\mathcal A_n(x)=O(n^{-2})$ uniformly. Thus the squared and cubed series
are normally convergent there and define meromorphic functions of $x$.
The identity theorem extends the formulas from the unit disc to the
connected punctured plane. Their removable values at zero agree with
\eqref{ct:eq:quadratic} and \eqref{ct:eq:cube0}.
\end{proof}

\subsection{A half-shift harmonic cube}
The half-integer recurrence gives
\[
 \mathcal A_n(\tfrac12)
   =8(\log2+H_n-H_{2n})-\frac4{2n+1}=4R_n.
\]
At $x=1/2$,
\[
 J=2-4\log2,\qquad J'=-4,\qquad J''=16-28\zeta(3).
\]
Equations~\eqref{ct:eq:A-square}--\eqref{ct:eq:A-cube} therefore give
\[
 \sum_n(2n+1)\mathcal A_n(\tfrac12)^2=4\zeta(2)-4,
 \qquad
 \sum_n(2n+1)\mathcal A_n(\tfrac12)^3=40\zeta(3)-64\log2.
\]
Division by $16$ and $64$, respectively, proves
\eqref{ct:eq:R-identities}.

There is also an elementary integral realization,
\begin{equation}\label{ct:eq:R-integral}
 R_n=\int_0^1t^{2n}\frac{1-t}{1+t}\dd t.
\end{equation}
Indeed $\int_0^1t^{2n}/(1+t)\dd t=\log2+H_n-H_{2n}$.
Thus one further exact integral consequence is
\begin{equation}\label{ct:eq:triple-rational}
 \int_{[0,1]^3}\frac{1+(xyz)^2}{(1-(xyz)^2)^2}
        \frac{1-x}{1+x}\frac{1-y}{1+y}\frac{1-z}{1+z}
           \dd x\dd y\dd z
       =\frac58\zeta(3)-\log2.
\end{equation}
Tonelli applies to the positive power series
$\sum_{n\ge0}(2n+1)(xyz)^{2n}$; its integral is finite by the already
proved cubic moment. There is no hidden endpoint finite part.

\subsection{A colliding-root Gamma example}
Taking $p=q=1/4$ gives $\sigma=3/16,\tau=1/32$. The cubic has a repeated
root, but the master identity is regular and yields
\begin{equation}\label{ct:eq:quartergamma}
 \sum_{n\ge0}\left[
 (n+\tfrac12)\left\{
 \frac{\Gamma(n+5/4)^2\Gamma(n+1/2)}
      {\Gamma(n+3/4)^2\Gamma(n+3/2)}-1\right\}
 -\frac1{32(n+1/2)}\right]
       =\frac{\pi+4\log2-7}{64}.
\end{equation}
This example illustrates why singular derivatives of individually labeled
roots should not be confused with singularities of a symmetric trace.

\section{Polylogarithmic representations and explicit primitives}
\label{ct:sec:polylogs}
For positive integers $s_1,\ldots,s_\ell$ and $|z|<1$, use
\begin{equation}\label{ct:eq:multiple-li}
 \Li_{s_1,\ldots,s_\ell}(z)
  =\sum_{m_1>\cdots>m_\ell\ge1}
       \frac{z^{m_1}}{m_1^{s_1}\cdots m_\ell^{s_\ell}}.
\end{equation}
Only the largest summation index carries $z$. This convention fixes the
primitive formulas and the endpoint $\Li_{\boldsymbol s}(1)=\zeta(\boldsymbol s)$.

\begin{proposition}[An ordinary polylogarithmic realization of the fourth power]
\label{ct:prop:polylog-quartic}
\begin{equation}\label{ct:eq:polylog-quartic}
 \sum_{n\ge0}(2n+1)T_3(n)^4
 =\zeta(10)+\zeta(5)^2+
     \frac1{360}\int_0^1\frac{(-\log t)^6\Li_3(t)}{1-t}\dd t.
\end{equation}
\end{proposition}
\begin{proof}
The geometric series and the defining series of $\Li_3$ give
\[
 \frac1{6!}\int_0^1\frac{(-\log t)^6\Li_3(t)}{1-t}\dd t
 =\sum_{m>k\ge1}\frac1{m^7k^3}=\zeta(7,3).
\]
All terms are positive. The integrand is integrable: near zero the factor
$\Li_3(t)$ is $O(t)$, and near one the sixth logarithmic power compensates
the denominator. Apply Corollary~\ref{ct:cor:quartic}.
\end{proof}

\begin{proposition}[An explicit primitive ladder]\label{ct:prop:primitive-li}
For $m\ge0$, $q\ge1$ an integer, and $0<z<1$,
\begin{equation}\label{ct:eq:primitive-li}
 \int_0^z\frac{(-\log t)^m\Li_q(t)}{1-t}\dd t
 =\sum_{j=0}^m\frac{m!}{(m-j)!}(-\log z)^{m-j}\Li_{j+1,q}(z).
\end{equation}
On $|z|<1$, the same identity holds with consistently chosen logarithms
and paths from zero. It can also be differentiated any fixed number of
times in the order $q$ when $\Li_q$ and the nested series are regarded as
entire functions of that order in the open disc.
\end{proposition}
\begin{proof}
The defining series give
\[
 \frac{\mathrm d}{\mathrm dz}\Li_{1,q}(z)=\frac{\Li_q(z)}{1-z},\qquad
 \frac{\mathrm d}{\mathrm dz}\Li_{r,q}(z)=\frac{\Li_{r-1,q}(z)}z\quad(r>1).
\]
Differentiating the finite sum on the right of
\eqref{ct:eq:primitive-li} telescopes; only the desired integrand remains.
Every term tends to zero as $z\to0$ along the chosen path because its nested
series starts at order $z^2$. This fixes the integration constant.
Local uniform geometric convergence, including every fixed order derivative,
justifies the last assertion. At $z=1$, the endpoint limit for $m\ge1$
is $m!\zeta(m+1,q)$ whenever the defining endpoint series converges.
\end{proof}
These integration rules are standard consequences of the iterated-series
calculus; their role here is to turn the depth-descent output into ordinary
polylogarithmic integrals and explicit antiderivatives.

\subsection{Generating functions for arbitrary powers}
For $s\ge3$, define
\[
 F_{s,d}(z)=\sum_{n\ge0}T_s(n)^d z^n,\qquad |z|<1.
\]
For $k\ge1$ put
\[
 V_{s,k}(z)=\sum_{\ell=1}^k\ 
       \sum_{m_1+\cdots+m_\ell=k\atop m_i\ge1}
       \frac{k!}{m_1!\cdots m_\ell!}
             \Li_{sm_1,\ldots,sm_\ell}(z),
 \qquad V_{s,0}(z)=1.
\]
Expanding the $k$th power of $H_n^{(s)}$ by the equal-index blocks gives
\begin{equation}\label{ct:eq:tail-gf}
 F_{s,d}(z)=\frac1{1-z}
      \sum_{k=0}^d(-1)^k\binom dk\zeta(s)^{d-k}V_{s,k}(z).
\end{equation}
There is no appeal to numerical multiple-polylogarithm evaluation in this
identity. The coefficient of $z^n$ is exactly
$(\zeta(s)-H_n^{(s)})^d$.

Writing $\theta=z\partial_z$,
\begin{equation}\label{ct:eq:weighted-gf}
 (2\theta+1)F_{s,d}(z)=\sum_{n\ge0}(2n+1)T_s(n)^d z^n.
\end{equation}
For $s=3,d\ge2$, absolute convergence permits $z\uparrow1$ on the
right. The cubic and quartic evaluations above are therefore explicit
endpoint cancellation identities for the multiple polylogarithms in
\eqref{ct:eq:tail-gf}. One must take the complete combination on the left;
its individual rationally weighted pieces need not have finite endpoint
limits.

\section{Algebraic Hurwitz traces and all-index Stieltjes identities}
\label{ct:sec:stieltjes}
We now use the same cubic roots, but allow a positive base point $h$ with
$h>\max_j|r_j|$. This condition ensures that every $h+r_j$ lies in the
right half-plane. Let
\begin{equation}\label{ct:eq:trace}
 \Ztr(s;h)=\sum_{j=1}^3\zeta(s,h+r_j)-3\zeta(s,h).
\end{equation}
The residues at $s=1$ cancel, so $\Ztr$ is entire in $s$.
The dependence on $\sigma,\tau$ is symmetric and holomorphic, including
at the cubic discriminant $4\sigma^3-27\tau^2=0$.

\begin{theorem}[Invariant spectral trace expansion]\label{ct:thm:trace}
For $h>\max_j|r_j|$,
\begin{equation}\label{ct:eq:trace-series}
 \Ztr(s;h)=
 \sum_{a,b\ge0\atop a+b>0}
 \frac{k}{a+b}\binom{a+b}{b}
       \frac{(s)_k}{k!}\zeta(s+k,h)\sigma^a\tau^b,
       \qquad k=2a+3b.
\end{equation}
The grouped series is locally normally convergent in $s$ and in parameter
sets on which all roots remain strictly inside $|r|<h$. An explicit
sufficient polydisc is $|\sigma|/h^2+|\tau|/h^3<1$.
Every apparent pole of an individual spectral factor at $s=1-k$ is
removable.
\end{theorem}
\begin{proof}
The ordinary Taylor expansion in the Hurwitz shift is
\[
 \zeta(s,h+r)=\sum_{k\ge0}\frac{(-1)^k(s)_k}{k!}\zeta(s+k,h)r^k.
\]
It follows initially from the absolutely convergent defining series for
$\Rea s>1$, and then by analytic continuation. The $k=0$ terms cancel
in the trace; the $k=1$ term vanishes because $\sum r_j=0$.
Substitute \eqref{ct:eq:newton}. Since $(-1)^k=(-1)^b$ when $k=2a+3b$,
the signs cancel and give \eqref{ct:eq:trace-series}.

For uniform convergence, choose a radius $\rho$ strictly between the largest
root modulus and $h$. The shift Taylor expansion is uniform on a larger
circle in the shift plane, on compact spectral sets after removal of the
common pole. Cauchy's estimate then gives a geometric majorant for the
root power sums. Each $k$ contains only finitely many invariant monomials;
on the stated smaller polydiscs their absolute expansions have the same
majorant. At $s=1-k$, the simple pole of $\zeta(s+k,h)$ is canceled by
the zero of $(s)_k$, which contains the factor $s+k-1$.
\end{proof}
\begin{remark}
The locally normal statement for an arbitrary root-bounded parameter set
uses the grouping by $k=2a+3b$. Absolute convergence of the unrestricted
double series is guaranteed on the explicit smaller polydisc. Keeping this
distinction avoids asserting absolute convergence merely from a favorable
choice of individually labeled roots.
\end{remark}

Define generalized Stieltjes constants by
\begin{equation}\label{ct:eq:stieltjes-convention}
 \zeta(1+u,h)=\frac1u+\sum_{m\ge0}\frac{(-1)^m}{m!}\gamma_m(h)u^m,
 \qquad \gamma_0(h)=-\psi(h),
\end{equation}
and set
\[
 \Delta_m(h)=\sum_{j=1}^3\gamma_m(h+r_j)-3\gamma_m(h).
\]
For $k\ge1$ let $E_\ell(k)$ be the elementary harmonic coefficients
\begin{equation}\label{ct:eq:E}
 \prod_{j=1}^k(1+u/j)=\sum_{\ell\ge0}\frac{E_\ell(k)}{\ell!}u^\ell.
\end{equation}
Thus $E_0=1$, $E_1=H_k$, $E_2=H_k^2-H_k^{(2)}$ and
$E_3=H_k^3-3H_kH_k^{(2)}+2H_k^{(3)}$; $E_\ell(k)=0$ for $\ell>k$.

\begin{corollary}[All generalized Stieltjes indices]\label{ct:cor:stieltjes}
In the same convergence domains,
\begin{align}
 \Delta_m(h)=(-1)^m\sum_{a,b\ge0\atop a+b>0}
   &\frac{k}{a+b}\binom{a+b}{b}\sigma^a\tau^b\notag\\
   &\times\sum_{\ell=0}^m\binom m\ell E_\ell(k)
                      \zeta^{(m-\ell)}(k+1,h),\qquad k=2a+3b.
 \label{ct:eq:stieltjes-trace}
\end{align}
The superscript on $\zeta$ denotes differentiation in its spectral argument.
\end{corollary}
\begin{proof}
Use $\Delta_m=(-1)^m\partial_s^m\Ztr(s;h)|_{s=1}$.
The factor $(1+u)_k/k!$ is exactly the product in \eqref{ct:eq:E}.
Leibniz's rule in \eqref{ct:eq:trace-series} gives the formula, and normal
convergence justifies termwise differentiation.
\end{proof}
For example, the first spectral derivative gives the explicit convergent
identity
\begin{equation}\label{ct:eq:gamma1}
 \Delta_1(h)=-\sum_{a,b\ge0\atop a+b>0}
 \frac{k}{a+b}\binom{a+b}{b}\sigma^a\tau^b
       \{\zeta'(k+1,h)+H_k\zeta(k+1,h)\}.
\end{equation}
This is a direct link between generalized Stieltjes constants, generalized
harmonic numbers, and derivatives of Hurwitz zeta in the cubic coordinates.

\subsection{The same trace is the right side of the Gamma master}
At $h=1$,
\begin{equation}\label{ct:eq:B-trace}
 B(\sigma,\tau)=\frac12-2\log2+\frac12\Ztr(1;1)
       =\frac12-2\log2+\frac12\Delta_0(1).
\end{equation}
Thus the right side of the nonlinear Gamma identity and the all-index
Stieltjes calculus have exactly the same invariant coefficients.
This does \emph{not} license differentiating the integer tail orders in the
nonlinear descent theorem: a separate continuously deformed summation
identity would be required for that stronger conclusion.

\subsection{A rational-shift check at every Stieltjes index}
Take $r=(1/2,0,-1/2)$, so $(\sigma,\tau)=(1/4,0)$, and $h=1$.
The Hurwitz recurrence and duplication give
\[
 \Ztr(s;1)=(2^{s+1}-4)\zeta(s)-2^s.
\]
If $L=\log2$ and $\gamma_j=\gamma_j(1)$, then
\begin{equation}\label{ct:eq:half-stieltjes}
 \Delta_m(1)=(-1)^m\left\{
 \frac{4L^{m+1}}{m+1}
 +4\sum_{j=0}^{m-1}\binom mj L^{m-j}(-1)^j\gamma_j
 -2L^m\right\},
\end{equation}
with an empty sum when $m=0$. In particular $\Delta_0(1)=4L-2$.
The pole is canceled \emph{before} Taylor expansion. Expanding the three
Hurwitz zeta functions separately at $s=1$ without canceling their common
pole is an avoidable source of incorrect constants.

\section{Gamma determinants and anchored antiderivatives}
\label{ct:sec:primitives}
Define
\begin{equation}\label{ct:eq:D}
 \Dtr(h;\sigma,\tau)=\frac{\prod_{j=1}^3\Gamma(h+r_j)}{\Gamma(h)^3},
\end{equation}
with its logarithm normalized to zero at $(\sigma,\tau)=(0,0)$.
The zeta derivative at zero and the cancellation of the three common
$\log(2\pi)/2$ terms imply
\begin{equation}\label{ct:eq:determinant}
 \log\Dtr(h;\sigma,\tau)=\Ztr'(0;h).
\end{equation}
At $h=1$, the elementary Gamma product gives
\begin{align}
 \Dtr(1;\sigma,\tau)
       &=\prod_{n\ge1}(1-\sigma/n^2-\tau/n^3)^{-1},\label{ct:eq:D-product}\\
 \log\Dtr(1;\sigma,\tau)
       &=\sum_{a,b\ge0\atop a+b>0}
           \frac1{a+b}\binom{a+b}{b}\zeta(2a+3b)\sigma^a\tau^b.
            \label{ct:eq:D-series}
\end{align}
Also $Q_n=\Dtr(n+1;\sigma,\tau)/\Dtr(n+1;\sigma,-\tau)$.
This identifies the odd-$\tau$ Gamma quotient that caused the depth drop.

\begin{proposition}[Derivative tower]\label{ct:prop:derivatives}
For every $r\ge0$,
\begin{equation}\label{ct:eq:trace-derivative}
 \partial_h^r\Ztr(s;h)=(-1)^r(s)_r\Ztr(s+r;h).
\end{equation}
Consequently, for $r\ge1$,
\begin{equation}\label{ct:eq:stieltjes-derivative}
 \partial_h^r\Delta_m(h)=(-1)^{m+r}r!
       \sum_{\ell=0}^m\binom m\ell E_\ell(r)
                     \Ztr^{(m-\ell)}(1+r;h).
\end{equation}
\end{proposition}
\begin{proof}
Apply the standard parameter derivative
$\partial_h\zeta(s,h)=-s\zeta(s+1,h)$ repeatedly to the finite trace.
Differentiate the resulting identity $m$ times in $s$ at $s=1$.
The Taylor coefficients of $(1+u)_r$ are $r!E_\ell(r)/\ell!$.
\end{proof}

We give every positive integer order of antiderivative with its integration
constants fixed. Choose real $h,h_0>\max_j|r_j|$; the formulas extend
holomorphically to an appropriate simply connected right-half-plane domain.
For $r\ge1$ set
\begin{align}
 N_r(u;h,h_0)={}&\Ztr(1-r+u;h)\notag\\
  &-\sum_{j=0}^{r-1}\frac{(h-h_0)^j}{j!}
         \left.\partial_t^j\Ztr(1-r+u;t)\right|_{t=h_0}.
 \label{ct:eq:N-primitive}
\end{align}

\begin{theorem}[All anchored Stieltjes primitives]\label{ct:thm:anchored}
The quotient $N_r(u;h,h_0)/(1-r+u)_r$ has a removable singularity at $u=0$.
Define
\begin{equation}\label{ct:eq:anchored}
 \mathcal I_{r,m}(h;h_0)=(-1)^{m+r}m![u^m]
                \frac{N_r(u;h,h_0)}{(1-r+u)_r}.
\end{equation}
Then
\begin{equation}\label{ct:eq:anchored-properties}
 \partial_h^r\mathcal I_{r,m}(h;h_0)=\Delta_m(h),\qquad
 \left.\partial_h^j\mathcal I_{r,m}(h;h_0)\right|_{h=h_0}=0
       \quad(0\le j<r).
\end{equation}
Equivalently,
\begin{equation}\label{ct:eq:anchored-integral}
 \mathcal I_{r,m}(h;h_0)
   =\frac1{(r-1)!}\int_{h_0}^{h}(h-t)^{r-1}\Delta_m(t)\dd t.
\end{equation}
\end{theorem}
\begin{proof}
The negative-integer identity
$\zeta(1-r,h)=-B_r(h)/r$ shows that $\Ztr(1-r;h)$ is a polynomial in $h$
of degree at most $r-2$: its leading degree-$r$ terms cancel between the
three shifts and the base, and its degree-$(r-1)$ coefficient is a multiple
of $\sum r_j=0$. For $r=1$ it vanishes identically. Its Taylor polynomial
of degree $r-1$ around $h_0$ is therefore exact, so $N_r(0;h,h_0)=0$.
The denominator has exactly one zero at $u=0$, with derivative
$(-1)^{r-1}(r-1)!$. Holomorphic division proves removability.

Differentiate \eqref{ct:eq:N-primitive} $r$ times in $h$. The subtracted
Taylor polynomial disappears and \eqref{ct:eq:trace-derivative} gives
\[
 \partial_h^r N_r(u;h,h_0)
       =(-1)^r(1-r+u)_r\Ztr(1+u;h).
\]
After division and coefficient extraction, the sign in
\eqref{ct:eq:anchored} gives $(-1)^m m![u^m]\Ztr(1+u;h)=\Delta_m(h)$.
The lower derivatives vanish at $h=h_0$ by construction, including after
the removable division. The usual uniqueness of the solution of the
$r$th derivative equation with these initial data gives
\eqref{ct:eq:anchored-integral}.
\end{proof}

For the first primitive this simplifies to
\begin{equation}\label{ct:eq:first-primitive}
 \mathcal I_{1,m}(h;h_0)=\frac{(-1)^{m+1}}{m+1}
       \{\Ztr^{(m+1)}(0;h)-\Ztr^{(m+1)}(0;h_0)\}.
\end{equation}
In particular,
\begin{equation}\label{ct:eq:gamma-primitive}
 \int_{h_0}^{h}\Delta_0(t)\dd t
       =\log\Dtr(h_0;\sigma,\tau)-\log\Dtr(h;\sigma,\tau).
\end{equation}
This is a Gamma antiderivative with no unspecified polynomial ambiguity.
The all-order formula treats that ambiguity explicitly rather than hiding
it in a symbol for a negative-order polygamma function.

\section{Lerch traces and polylogarithmic order derivatives}
\label{ct:sec:lerch}
Let $|z|<1$ and
\[
 \Phi(z,s,h)=\sum_{n\ge0}\frac{z^n}{(n+h)^s},\qquad
 \Ltr(z,s;h)=\sum_{j=1}^3\Phi(z,s,h+r_j)-3\Phi(z,s,h).
\]
All powers use the logarithm in the right half-plane. The geometric factor
makes these functions entire in $s$. Exactly the proof of
Theorem~\ref{ct:thm:trace} gives
\begin{equation}\label{ct:eq:lerch-trace}
 \Ltr(z,s;h)=\sum_{a,b\ge0\atop a+b>0}
 \frac{k}{a+b}\binom{a+b}{b}\frac{(s)_k}{k!}
        \Phi(z,s+k,h)\sigma^a\tau^b,\qquad k=2a+3b.
\end{equation}
At $h=1$, $z\Phi(z,s,1)=\Li_s(z)$, so
\begin{equation}\label{ct:eq:polylog-trace}
 z\sum_{j=1}^3\Phi(z,s,1+r_j)-3\Li_s(z)
 =\sum_{a,b\ge0\atop a+b>0}
   \frac{k}{a+b}\binom{a+b}{b}\frac{(s)_k}{k!}
      \Li_{s+k}(z)\sigma^a\tau^b.
\end{equation}
Every fixed spectral derivative may be taken termwise. For example,
\begin{align}
 \left.\partial_s\left\{
 z\sum_j\Phi(z,s,1+r_j)-3\Li_s(z)\right\}\right|_{s=1}
   =\sum_{a,b\ge0\atop a+b>0}
       &\frac{k}{a+b}\binom{a+b}{b}\sigma^a\tau^b\notag\\
       &\times\left\{\left.\partial_v\Li_v(z)\right|_{v=k+1}
                             +H_k\Li_{k+1}(z)\right\}.
 \label{ct:eq:polylog-order-derivative}
\end{align}
At higher derivative orders, the bracket is replaced by
$\sum_{\ell=0}^m\binom m\ell E_\ell(k)
\partial_v^{m-\ell}\Li_v(z)|_{v=k+1}$.

\subsection{The endpoint and why it is controlled}
For fixed shifts with zero sum,
\[
 \sum_j(n+h+r_j)^{-s}-3(n+h)^{-s}=O(n^{-\Rea s-2})
\]
locally uniformly in $s$, including fixed spectral derivatives with an
additional logarithmic factor. Thus the trace defining series at $z=1$
converges ordinarily for $\Rea s>-1$. On this half-plane Abel passage from
\eqref{ct:eq:lerch-trace} recovers \eqref{ct:eq:trace-series}. In particular
all spectral derivatives at $s=1$ have ordinary convergent endpoint traces.
For $\Rea s\le-1$, analytic continuation is available but ordinary endpoint
convergence is not asserted.

\subsection{A finite differential system in the argument}
The traces with powers of the roots also remain finite-dimensional. Put
\[
 F_j(z,s;h)=\sum_{i=1}^3r_i^j\Phi(z,s,h+r_i),\qquad j=0,1,2.
\]
Since $\theta\Phi(z,s,a)=\Phi(z,s-1,a)-a\Phi(z,s,a)$,
\begin{align}
 \theta F_0&=F_0(z,s-1;h)-hF_0-F_1,\notag\\
 \theta F_1&=F_1(z,s-1;h)-hF_1-F_2,\label{ct:eq:lerch-system}\\
 \theta F_2&=F_2(z,s-1;h)-hF_2-\sigma F_1+\tau F_0.\notag
\end{align}
The last line uses $r_i^3=\sigma r_i-\tau$.
This supplies argument derivatives without differentiating branch choices
for individual cubic roots.

Anchored $h$-antiderivatives of the order jets of $\Ltr$ are obtained by
replacing $\Ztr$ with $\Ltr$ in \eqref{ct:eq:N-primitive}--\eqref{ct:eq:anchored}.
The same removal proof works: at nonpositive integer order the Lerch series
is a polynomial in $h$ with coefficients rational in $z$, and the trace
cancels its leading terms. At $z=1$ one must use the Hurwitz trace version
or a justified endpoint limit rather than substituting into an individual
negative-order Lerch series.

\section{Audit findings and claims that must not be made}
\label{ct:sec:audit}
\subsection{The inspected mixed-resonance proof}
No new mathematical error was established in the inspected canonical
mixed-resonance chapter. Its centered $x_n=n+a/2$ normalization, the first
resonant subtraction, and the sign of the cubic counterterm at $a=1$ agree with
the independent calculation in Section~\ref{ct:sec:master-proof}.
The explicit Bernoulli cancellations are also checked by exact symbolic
arithmetic in the supplied code. This is a targeted audit, not an assertion
that every statement in the manuscript or incoming archives has been checked.

The necessary integration safeguards are mathematical rather than cosmetic.
The $n=0$ term must be included in the present moments; omitting it changes
an unweighted tail-power identity by the corresponding product of zeta
values. The weight here is $2n+1$, not $n+1/2$. The roots of the cubic have
product $-\tau$, whereas the odd parameter of the Gamma quotient is $+\tau$.
At $x=0$ the complete right side of the digamma formulas must be evaluated
by a removable limit. At higher Stieltjes index, the spectral pole must be
canceled before taking Taylor coefficients.

\subsection{A false extrapolation, not an attributed repository conjecture}
The cases $d=2,3$ might suggest
\begin{equation}\label{ct:eq:false-extrapolation}
 \sum_n(2n+1)T_3(n)^d\stackrel{?}{=}
       (2-2^{1-d})\zeta(3d-2).
\end{equation}
This is false already at $d=4$. We do not attribute
\eqref{ct:eq:false-extrapolation} to any existing source. A rational
countercheck is enough: positivity gives a lower bound
$\bigl(\sum_{j=1}^{4}j^{-3}\bigr)^4$ for the left side, while
\[
 \frac{15}{8}\zeta(10)
   <\frac{15}{8}\left(1+2^{-10}+\frac1{9\cdot2^9}\right)
   <\left(\sum_{j=1}^{4}j^{-3}\right)^4.
\]
The last inequality is an exact comparison of rational numbers. The correct
fourth-power formula is \eqref{ct:eq:quartic}. This small rejection is
included to prevent a misleading simplification, not because inequalities
are a research target of this article.

\subsection{The Gaussian conjectures are not solved here}
The canonical $S_4$ identity already has an exact octahedral certificate;
the inspected discovery chapter distinguishes a rigorously rejected old
$S_8$ vector from the surviving revised candidate \cite{ct:repo-status}.
The present descent concerns ordinary odd zeta tails. We have not proved a
bridge from the mixed-color Gaussian $S_6$ or revised $S_8$ residual to this
invariant family. Their status is unchanged by this work. Nor does the
appearance of $\zeta(7,3)$ prove any arithmetic irreducibility.

\section{Reproducibility and manuscript integration}
\label{ct:sec:verification}
The analytic proofs above are independent of the numerical computations.
The package includes one self-contained \LaTeX{} source, its compiled PDF,
verification code, machine-readable results, a coefficient catalog, an
integration note, and a compact theorem insert.

\subsection{Finite exact checks}
The exact suite makes \textbf{596 assertions}. It compares two independent
constructions of the invariant exponential coefficients through
$0\le a\le4$, $0\le d\le8$; checks their tail weight and the two-factor
parity descent; compares the Newton recurrence with the closed formula
through degree $30$; checks the harmonic/Pochhammer coefficients through
$k=20$ and order five; verifies finite minimum-square sums against ordered
partitions using rational arithmetic; checks the four printed low-degree
polynomials and the Bernoulli constants; and verifies the negative-order
Taylor cancellation for the first eight primitive orders.

The catalog records the $35$ nonlinear pairs $0\le a\le4$, $2\le d\le8$.
These tests verify finite algebra relative to the formulas stated in the
article. They are not a proof-assistant certification of Gamma continuation,
Tonelli's theorem, or normal convergence.

\subsection{Independent numerical diagnostics}
The numerical program uses defining Hurwitz tails for the quadratic,
cubic, quartic and quintic identities. The depth-two coordinate in the
quartic formula is evaluated independently by its polylogarithmic integral.
The digamma formulas are checked at real points on both sides of a pole
and at a nonreal point, always avoiding the pole itself. Further checks
compare the full subtracted Gamma master with its digamma value, the
invariant Stieltjes expansion with real-argument Stieltjes constants at a
repeated-root configuration, anchored primitives with ordinary quadrature,
and the finite polylogarithmic primitive with its defining integral.

All \textbf{31 numerical checks} passed at \textbf{50 decimal working digits}.
The output files record working precision, residual, tolerance and pass
status for every test actually completed. Numerical tolerances are
floating-point diagnostic thresholds, not rigorous error enclosures.
In particular, no integer-relation search is used to infer an identity.

\subsection{Rebuilding and proposed placement}
From the package directory, run
\par\medskip\noindent\begin{minipage}{\linewidth}
\begin{verbatim}
python verification/verify.py --exact-only
python verification/verify.py --numeric-only --dps 50
pdflatex -interaction=nonstopmode -halt-on-error article.tex
pdflatex -interaction=nonstopmode -halt-on-error article.tex
\end{verbatim}
\end{minipage}
\par\medskip
The proposed report home is the continuation directory
\begin{quote}\small
\src{Analysis/Polylogarithms/docs/reports/stieltjes-harmonic-resonance/}\\
\src{continuations/cubic-invariant-tail-descent/}.
\end{quote}
The main insertion follows the mixed-resonance chapter. The
trace and primitive results can instead
be cross-referenced from the Hurwitz differentiation and integration
chapters. All mathematical labels in this article use the prefix
\src{ct:}; the supplied compact insert uses \src{cti:}. Proposed intake and
integration are not claimed to have occurred.

\section{Further research questions}
\label{ct:sec:future}
\paragraph{1. An elementary certificate for the master identity.}
Can \eqref{ct:eq:master} be proved by an explicit finite-difference or
contour certificate that avoids the meromorphic continuation of the
Rogers--Dougall sum? The cubic product ratio
$Q_{n-1}/Q_n=(n^3-\sigma n+\tau)/(n^3-\sigma n-\tau)$ is unusually simple.
A successful certificate must recover the nonzero constant
$1/2-2\log2$, not just its nonlinear coefficients. Such a proof could also
provide a shorter analytic foundation for formalization.

\paragraph{2. Full polarization of the symmetric-digamma cube.}
Theorem~\ref{ct:thm:digamma} evaluates the diagonal sum
$\sum_n(2n+1)\mathcal A_n(x)^3$. Determine a closed formula for
$\sum_n(2n+1)\mathcal A_n(x)\mathcal A_n(y)\mathcal A_n(z)$ at independent
arguments. Derivatives on the diagonal and the symmetric convolutions
proved here do not uniquely determine this three-variable function.
A divided-difference formula with controlled collision limits would be
particularly useful.

\paragraph{3. The fifth-power normal form.}
Equation~\eqref{ct:eq:339} provides a small, explicit weight-thirteen
starting basket. Reduce its particular combination further by exact
shuffle, stuffle and duality relations, keeping a rational certificate.
The target is a shorter proven expression for the fifth moment, not a
numerical claim that a candidate list is linearly independent.

\paragraph{4. Beyond three roots.}
For four or more roots with zero sum, the odd part of a Gamma determinant
still has a parity filtration, but there are additional independent
invariants. Determine which higher hypergeometric summations supply a
normally convergent weighted trace identity strong enough to isolate the
leading tail moments. Parity of an exponential alone does not supply the
missing summation theorem.

\paragraph{5. Colored analogues and the Gaussian obstruction.}
Develop a roots-of-unity analogue of the cubic product in which the tail
indices are restricted to congruence classes or carry colors. A meaningful
first target is an exact formula for a colored cubic tail moment whose
ordered-partition image contains the mixed-color Gaussian depth-two
coordinates. Only a proved transformation of that kind could connect
this descent to $S_6$ or the revised $S_8$ conjecture.

\paragraph{6. A genuine continuous-order nonlinear deformation.}
The Lerch trace permits arbitrary order derivatives, whereas the
Rogers--Dougall tail descent uses the fixed odd integer grading. Construct
a nonlinear master identity with a free spectral order that interpolates
the integer moments and whose derivatives produce weighted logarithmic
tails. One must not obtain such a deformation by differentiating a family
known only at isolated integer orders.

\paragraph{7. Formal verification with analytic hypotheses.}
The recurrence, Newton algebra, ordered-partition map and primitive
removal are finite algebraic targets. A useful next step is to formalize
these first, with the normally convergent master identity as an explicit
hypothesis, and then formalize the Gamma-continuation and interchange
lemmas. The supplied computations already separate those two layers.

\section*{Conclusion}
The cubic invariant coordinates expose a triangular feature of an existing
Gamma resonance: nonlinear coefficients of high tail degree are controlled
by lower tail degree, with a two-step parity gap. This yields an exact
all-degree descent, particularly simple cubic and quartic evaluations,
and meromorphic formulas for symmetric digamma powers. The same invariants
organize Hurwitz, Stieltjes, Gamma and Lerch traces without introducing
root-branch singularities, and allow every integration constant in the
primitive tower to be specified. The result is a concrete continuation of
the exact-identity program, with clearly separated classical inputs,
proved consequences, numerical diagnostics and remaining conjectures.

\clearpage
\begin{thebibliography}{9}
\raggedright
\bibitem{ct:repo-mixed}
ProveIt Contributors,
\emph{Mixed Gamma derivatives and exact harmonic tail identities},
canonical manuscript chapter
\src{Analysis/Polylogarithms/docs/manuscript/chapters/07-mixed-dougall-identities.tex},
source anchor \src{87c7db629552afbc2e803001526debe3430a31b0}, inspected 10 October 2026.
Repository: \url{https://github.com/VladimirReshetnikov/ProveIt}.

\bibitem{ct:repo-harmonic}
ProveIt Contributors,
\emph{Harmonic Gamma coefficient identities},
canonical chapter \src{07-harmonic-gamma-coefficients.tex}, same repository
and source anchor as \cite{ct:repo-mixed}.

\bibitem{ct:repo-status}
ProveIt Contributors,
canonical manuscript \src{README.md}, \src{chapters/04-S4-proof.tex}
and \src{chapters/10-discovery.tex}; incoming \src{README.md} and directory
inventory, inspected 10 October 2026. Same repository as \cite{ct:repo-mixed}.
These are source-status references, not assertions of a complete archive audit.

\bibitem{ct:dlmf16}
NIST Digital Library of Mathematical Functions,
\S16.4, especially Eq.~16.4.9, Rogers--Dougall very well-poised sum.
\url{https://dlmf.nist.gov/16.4.E9}. Consulted 10 October 2026.

\bibitem{ct:dlmf5}
NIST Digital Library of Mathematical Functions,
\S5.8, infinite products, and \S5.11, asymptotic expansions of Gamma
functions and ratios.
\url{https://dlmf.nist.gov/5.8}; \url{https://dlmf.nist.gov/5.11}.
Consulted 10 October 2026.

\bibitem{ct:dlmf25}
NIST Digital Library of Mathematical Functions,
\S25.11, Hurwitz zeta function: shift recurrence, parameter derivatives,
negative-integer values and $\zeta'(0,h)$; \S25.14, Lerch's transcendent.
\url{https://dlmf.nist.gov/25.11}; \url{https://dlmf.nist.gov/25.14}.
Consulted 10 October 2026.

\bibitem{ct:hoffman}
Michael E. Hoffman,
\emph{Sums of Products of Riemann Zeta Tails},
arXiv:1602.03399, 2016. In particular \S2 for the Euler depth-two identities
and \S\S3--4 for products of tails.
\url{https://arxiv.org/abs/1602.03399}.
\end{thebibliography}
\end{document}
